%%%%%%%%%%%%%%%%%%%%%%%%%%%%%%%%%%%%%%%%%%%%%%%%%%%%%%%%%%%%%%%%%%%%%%%
\begin{comment}
\textblockcolour{Lavender}
\begin{textblock}{12.5}(0.5, 1)
    \noindent
    ピンクは重要項目、黄色はより良い論文にするための項目、緑は任意項目、青は補足説明、先頭の数字はチェックリストの番号に対応する
\end{textblock}

\textblockcolour{pink}
\begin{textblock}{5}(0.5, 3)
    【8】紙印刷の頃にあった

    背表紙を作らない
\end{textblock}

\textblockcolour{pink}
\begin{textblock}{4}(0.5, 8)
    【6,10】タイトル
\end{textblock}

\textblockcolour{pink}
\begin{textblock}{4}(0.5, 10)
    【6,10】指導教員名
\end{textblock}

\textblockcolour{pink}
\begin{textblock}{4}(0.5, 16.5)
    【6,10】学部・研究室
\end{textblock}

\textblockcolour{pink}
\begin{textblock}{4}(0.5, 19)
    \noindent
    【6,11】学籍番号・氏名
\end{textblock}

\textblockcolour{pink}
\begin{textblock}{4}(16, 1)
    【1】1枚目は「表紙」
\end{textblock}

\textblockcolour{lime}
\begin{textblock}{4}(14, 11)
    ↑職位の有無は任意
\end{textblock}



\begin{comment19}
    \textblockcolour{pink}
    \begin{textblock}{11}(8, 15)
        ↓19年度以前の入学生は学部の後に研究室名を記載する
    \end{textblock}
\end{comment19}

\begin{comment20}
    \textblockcolour{pink}
    \begin{textblock}{8}(8, 14)
        ↓20年度以降の入学生は学部の後に専攻名、

        次の行に研究室名を記載する
    \end{textblock}
\end{comment20}

\textblockcolour{pink}
\begin{textblock}{6}(10, 19)
    【12】学籍番号のA,B,Cは大文字
\end{textblock}

\textblockcolour{pink}
\begin{textblock}{4}(12, 24)
    【13】↑2026年度
\end{textblock}

\begin{textblock}{7}(11, 28)
    ←外表紙にページ番号をつけない
\end{textblock}
\end{comment}
%%%%%%%%%%%%%%%%%%%%%%%%%%%%%%%%%%%%%%%%%%%%%%%%%%%%%%%%%%%%%%%%%%%%%%%